\section{Проверка свободных энергий в малых ячейках}
\label{sec:smallcells}

Малый размер ячейки позволил выполнить полный перебор химических заполнений и
на одних и тех же конфигурациях измерить два эффекта. Прямые фононы показывают,
насколько кластерное разложение C46 воспроизводит гармонические свободные
энергии MTP. Переход от гармонического гамильтониана к полному MTP при
классической статистике ядер измеряет вклад ангармонизма. Все расчёты этого
раздела относятся к $T=1000$~К, фиксированным $V/N=17{,}28833$~\AA$^3$/атом
и $c/a=4{,}6046/2{,}9446$. Здесь вычислена свободная энергия Гельмгольца.

\subsection{Гармоническая опора для каждого химического заполнения}

Для каждого симметрийно независимого заполнения $\sigma$ фиксировали векторы
ячейки и сорта атомов, релаксировали внутренние координаты и записывали
энергию минимума $U_{\min,\sigma}$. Гессиан $H_\sigma$ получали центральными
разностями сил того же MTP. При смещении координаты $j$ на $\pm h$ элемент
гессиана равен $-[F_i(+h)-F_i(-h)]/(2h)$. Матрицу симметризовали. На трёх
заполнениях восьмиатомной ячейки уменьшение $h$ с $0{,}005$ до
$0{,}0025$~\AA\ изменило классическую гармоническую свободную энергию максимум
на $0{,}0029$~мэВ/атом. Для перевзвешивания использовали меньший шаг.

Общий перенос ячейки даёт три нулевые моды. Ортонормированный базис $B$
выбирали в подпространстве смещений, для которых $\sum_i\mathbf u_i=0$.
Собственные числа $\kappa_j$ матрицы $B^{\mathsf T}H_\sigma B$ должны быть
положительны. Если $S$ содержит соответствующие собственные векторы, то
тепловые смещения относительно релаксированных координат выражаются через
безразмерный вектор $\boldsymbol\xi$,
\begin{equation}
 \mathbf u=B S\,\mathrm{diag}\!\left(\sqrt{k_{\rm B}T/\kappa_j}\right)
  \boldsymbol\xi,\qquad
 U_{\mathrm h}=U_{\min,\sigma}+\tfrac12 k_{\rm B}T
                    \lVert\boldsymbol\xi\rVert^2 .
\end{equation}
В неограниченном гармоническом ансамбле компоненты $\boldsymbol\xi$
независимы и имеют стандартное нормальное распределение с совместной
плотностью $\phi$. Его размерность $d=3N-3$ равна 21 или 45. Гауссовский
интеграл задаёт опорную свободную энергию $F_{\sigma,\mathrm h}^{\infty}$.
Для сравнений при $N=16$ достаточно её конфигурационной части
$U_{\min,\sigma}+\tfrac12k_{\rm B}T\sum_j
\ln[\kappa_j\ell_0^2/(2\pi k_{\rm B}T)]$ с общей для всех заполнений единицей
длины $\ell_0$. В сравнении при $N=8$ и одинаковом составе использовали
классическую модовую сумму $k_{\rm B}T\sum'_j\ln[\hbar\omega_j/(k_{\rm B}T)]$
в дополнение к $U_{\min,\sigma}$, где штрих исключает общие переносы.
В разностях между заполнениями этого состава общий массовый вклад сокращается.

\subsection{Термодинамический интеграл и вычисление средних}

При $\beta=(k_{\rm B}T)^{-1}$ и фиксированном заполнении вводим
$\Delta U=U_{\mathrm{MTP}}-U_{\mathrm h}$ и
$U_\lambda=U_{\mathrm h}+\lambda\Delta U$. Ограничение на координаты $D$
выбирается одинаковым на всём пути. Оно отсчитывается от идеальных узлов
ГПУ, тогда как гармоническое смещение отсчитывается от релаксированного
минимума. Интегрирование по параметру связи предложил
Кирквуд~\cite{kirkwood1935}. Дифференцирование
$F_{\sigma,\lambda}^D=-k_{\rm B}T\ln Z_{\sigma,\lambda}^D$ даёт
\begin{equation}
 I_\lambda=\frac{\partial F_{\sigma,\lambda}^D}{\partial\lambda}
  =\langle\Delta U\rangle_{\sigma,\lambda,D},\qquad
 F_{\sigma,\mathrm{MTP}}^D-F_{\sigma,\mathrm h}^D
  =\int_0^1 I_\lambda\,\mathrm d\lambda .
\end{equation}
На одном конце пути восстанавливаются колебания в квадратичной яме, на другом
каждое смещение оценивается полной энергией MTP. $I_\lambda$ имеет
размерность энергии ячейки. Деление на $N$ выполняется при представлении
результата на атом.

В пилотах промежуточные ансамбли восстанавливали по общей выборке независимых
пробных смещений. Основная выборка поровну содержала векторы из $\phi$ и из
изотропной гауссианы с шириной $1{,}35$ в тех же координатах. Для каждой
пробы $s$ вычисляли полную плотность смеси $q(\boldsymbol\xi_s)$, значение
$\Delta U_s$ внутри $D$ и индикатор $\mathbf 1_D$. Массив энергий использовали
при всех значениях $\lambda$. При общем числе проб $L$ веса и оценки имеют вид
\begin{align}
 w_{s,\lambda}
 &=\mathbf 1_D(\boldsymbol\xi_s)\,
   \frac{\phi(\boldsymbol\xi_s)}{q(\boldsymbol\xi_s)}
                 \exp(-\beta\lambda\Delta U_s),&
 A_\lambda&=\frac{Z_{\sigma,\lambda}^D}{Z_{\sigma,\mathrm h}^{\infty}}
          \simeq \frac1L\sum_{s=1}^L w_{s,\lambda},\\
 \bar w_{s,\lambda}&=\frac{w_{s,\lambda}}{\sum_t w_{t,\lambda}},&
 I_\lambda&\simeq\sum_s\bar w_{s,\lambda}\Delta U_s .
\end{align}
Суммы экспонент вычисляли через логарифмы весов, за границей задавали
$\ln w=-\infty$. В $L$ входят также отвергнутые пробы. При $\lambda=0$
получаем $A_0=P_{\mathrm h}^\infty(D)$, вероятность бассейна в
неограниченном гармоническом ансамбле. Поэтому свободные энергии
ограниченных гармоники и MTP восстанавливаются по отдельности,
\begin{align}
 F_{\sigma,\mathrm h}^D&=F_{\sigma,\mathrm h}^\infty-k_{\rm B}T\ln A_0,&
 F_{\sigma,\mathrm{MTP}}^D&=F_{\sigma,\mathrm h}^\infty-k_{\rm B}T\ln A_1 .
\end{align}
Их разность равна $-k_{\rm B}T\ln(A_1/A_0)$. Это формула экспоненциального
возмущения свободной энергии (FEP) Цванцига~\cite{zwanzig1954} в
ограниченном гармоническом ансамбле. Результат по концам пути
сравнивали с интегралом от $I_\lambda$. Отношение учитывает одинаковую
стенку и в числителе, и в знаменателе. Для каждого из ограничений
$W<0{,}35$ и $W<0{,}40$ табличные свободные энергии вычислены именно по
концевым весам.

Первоначальную проверку при $N=8$ проводили для всей ячейки Вороного
($W<1$). Для каждого из трёх классов сгенерировали восемь независимых
блоков по 2048 проб. Интеграл вычисляли трапециями по 65 равноотстоящим
узлам, включая оба конца. Переход к 17 узлам менял результат не более чем
на $0{,}0012$~мэВ/атом. Расхождение между 65 узлами и концевой оценкой
составило менее $0{,}000074$~мэВ/атом. Поскольку оценки используют одни и
те же смещения, это проверка квадратуры. Статистическую концентрацию весов
характеризует эффективное число проб $N_{\mathrm{eff}}=1/\sum_s\bar
w_{s,\lambda}^2$. Минимумы по пути составили примерно 7477 для
D0$_{19}$, 4397 для межслойной и 6003 для базальной пары из 16\,384 проб на
класс. Максимальный вес одной пробы на конце пути достигал $0{,}50\%$.

Независимой проверкой $I_1$ служили цепи Метрополиса в том же бассейне при
$\lambda=1$. Ход $\boldsymbol\xi'=\rho\boldsymbol\xi+
\sqrt{1-\rho^2}\boldsymbol\eta$, $\boldsymbol\eta\sim\phi$, сохраняет
гармоническую гауссиану. Для проб внутри $D$ вероятность принятия равна
$\min\{1,\exp[-\beta(\Delta U'-\Delta U)]\}$. Для каждого класса запускали
48 цепей при $\rho=0{,}7$, отбрасывали 128 начальных шагов и анализировали
следующие 384 шага каждой цепи. Половина цепей стартовала из минимума,
остальные из различных смещённых состояний. Прямое среднее
и перевзвешивание различались максимум на $1{,}60$ объединённой
стандартной ошибки. Обе численные проверки относятся к посещённым
конфигурациям. Независимый поиск более редких состояний описан ниже.

Чтобы получить каноническую свободную энергию при заданном числе атомов Al
$k$, суммировали статистические веса по различным заполнениям,
\begin{equation}
 F_k^D=-k_{\mathrm B}T\ln
   \sum_{\sigma\,:\,N_{\mathrm{Al}}=k}g_\sigma
     \exp[-\beta F_\sigma^D] .
\end{equation}
Здесь суммирование идёт по одному представителю класса. Кратность
$g_\sigma$ включает все его симметрийно эквивалентные заполнения,
включая различные ориентации упорядочения. После этой суммы дополнительный
множитель для ориентаций привёл бы к повторному учёту.

\subsection{Короткие расстояния и выбор кристаллического бассейна}

В независимой проверке межслойного заполнения $\mathrm{Ti}_6\mathrm{Al}_2$
сгенерировали 65\,536 пробных смещений из смеси с равными долями
стандартного нормального распределения и изотропного многомерного распределения
Стьюдента с пятью степенями свободы и масштабом $1{,}20$ в координатах
$\boldsymbol\xi$. Компонента Стьюдента чаще предлагает состояния с большими
смещениями. В одном состоянии два атома Ti
сблизились до $0{,}636816$~\AA\ при сохранении каждого атома в своей исходной
ячейке Вороного. Полная энергия восьми атомов упала с $-56{,}069$ до
$-128\,392{,}623$~эВ. Почти весь статистический вес пришёлся на один снимок,
и эффективное число проб снизилось практически до единицы.
Повторный расчёт MLIP и периодическое размножение до 64 атомов воспроизвели
аномалию. Расхождение энергий на атом при размножении составило
$2{,}85\cdot10^{-9}$~эВ. На прямолинейном пути из 101 точки между релаксированной
структурой и этим снимком энергия сначала поднималась на $42{,}20$~эВ/ячейку
над энергией минимума. Высота относится к конкретному пути и объясняет, почему
локальные траектории способны долго удерживаться вблизи кристалла.

Кристаллический бассейн задали через смещение каждого атома от идеального
ГПУ-узла после удаления общего переноса. Для каждой грани его ячейки
Вигнера--Зейтца вычисляли нормальную проекцию смещения в долях расстояния
от узла до грани. Максимум этих безразмерных величин по всем атомам и
граням обозначили $W$. При $W<0{,}40$ минимально возможное расстояние в
этой ячейке ограничено снизу величиной
$(1-0{,}40)\cdot 2{,}861976 = 1{,}71719$~\AA. Параметр радиального базиса
MTP равен $1{,}68052$~\AA. Опубликованный файл потенциала допускает
вычисление энергии и на меньших расстояниях. Геометрическое ограничение
поэтому входит в определение измеренного кристаллического равновесия.

\subsection{Предпочтение порядка \texorpdfstring{D0$_{19}$}{D0₁₉} в ячейке из восьми атомов}

При $N=8$ и $k=2$ имеется $\binom82=28$ заполнений. Группировка по симметрии
даёт 4 варианта D0$_{19}$, 12 вариантов с базальной парой Al и 12 вариантов
с межслойной парой другого типа. В табл.~\ref{tab:small8} приведена энергия
одного представителя каждого класса относительно D0$_{19}$. Кратности входят
в статистическую сумму только при последующем объединении заполнений.
В двух последних строках таблицы для D0$_{19}$ и базальной пары использована
смесь двух гауссиан. Для межслойной пары использованы пробы независимой серии
со Стьюдентом, отфильтрованные по тому же $W<0{,}40$. Если для всех трёх
классов брать исходную гауссовскую смесь, ангармоническая поправка к цене
межслойной пары составляет $2{,}28\pm0{,}18$~мэВ/атом.

\begin{table}[htbp]
\centering\small
\caption{Энергетическая цена двух заполнений $\mathrm{Ti}_6\mathrm{Al}_2$
относительно D0$_{19}$, мэВ/атом. Знак $\pm$ здесь и в
табл.~\ref{tab:small16} обозначает одну стандартную ошибку.}
\label{tab:small8}
\begin{tabular}{@{}lrr@{}}
\hline
Модель & Базальная пара & Межслойная пара\\
\hline
C46, статическая энергия & $78{,}60$ & $110{,}61$\\
Прямой MTP, статическая энергия & $84{,}10$ & $112{,}92$\\
C46, квантовая гармоника & $61{,}01$ & $85{,}66$\\
Прямой MTP, квантовая гармоника & $65{,}43$ & $86{,}44$\\
Классическая гармоника, $W<0{,}40$ & $65{,}55\pm0{,}05$ & $87{,}78\pm0{,}05$\\
Классический MTP, $W<0{,}40$ & $68{,}07\pm0{,}16$ & $89{,}98\pm0{,}10$\\
\hline
\end{tabular}
\end{table}

Для прямой квантовой гармоники силовые постоянные строили в суперячейках
$5\times5\times6$ исходных восьмиатомных ячеек (1200 атомов), фононную
сумму вычисляли на сетке $16^3$ в зоне Бриллюэна. При общей гармонической
модели C46 занижает цену базальной пары на $4{,}42$~мэВ/атом, межслойной
на $0{,}78$~мэВ/атом. В восьмиатомной классической ячейке ангармонизм
повышает эти цены соответственно на $2{,}52\pm0{,}13$ и
$2{,}20\pm0{,}08$~мэВ/атом по согласованным разностям. Независимая широкая
выборка для всех трёх заполнений дала $2{,}08\pm0{,}17$ и
$1{,}60\pm0{,}65$~мэВ/атом. Обе серии дают дополнительную стабилизацию порядка.
Для базальной пары расхождение оценок составляет около двух объединённых
стандартных ошибок. На рис.~\ref{fig:ti-integrand} показано, как
получены интегральные поправки для трёх заполнений. Квантовые результаты включают колебания во всей зоне
Бриллюэна, классические относятся к модам малой периодической ячейки.
Суммирование двух измеренных сдвигов потребовало бы согласования этих условий.

\subsection{Нелинейность свободной энергии по составу в ячейке на 16 атомов}

Для $N=16$ и $k=2,3,4$ атомные доли Al равны $12{,}5$, $18{,}75$ и
$25$~ат.\,\%. Перечислены соответственно 120, 560 и 1820 заполнений. После
учёта симметрии осталось 5, 9 и 29 классов. Для каждого из 43 представителей
получили устойчивый гессиан по 45 внутренним степеням свободы. Меру
нелинейности определили как
\[
 \delta f_{16}=\frac{F_3-(F_2+F_4)/2}{16}.
\]
Положительная $\delta f_{16}$ означает, что при среднем составе свободная
энергия лежит выше прямой между крайними составами. Комбинация сокращает
общие постоянные и линейные по $k$ члены, включая классический вклад масс
от интегрирования по импульсам. У статических гамильтонианов после
суммирования химических заполнений $\delta f_{16}$ составила
$3{,}015$~мэВ/атом для C46 и $3{,}706$~мэВ/атом для прямого MTP.

\begin{table}[htbp]
\centering\small
\caption{Значение $\delta f_{16}$ и его ангармоническая поправка, мэВ/атом.
Столбцы представляют независимые серии. Во второй использована смесь
гауссовского распределения и распределения Стьюдента с пятью степенями свободы.}
\label{tab:small16}
\begin{tabular}{@{}lrr@{}}
\hline
Величина & Смесь гауссиан & Смесь со Стьюдентом\\
\hline
Гармоника, $W<0{,}40$ & $1{,}080\pm0{,}019$ & $1{,}077\pm0{,}009$\\
Полный MTP, $W<0{,}40$ & $1{,}374\pm0{,}085$ & $1{,}383\pm0{,}036$\\
Ангармоническая поправка, $W<0{,}40$ & $+0{,}294\pm0{,}071$ & $+0{,}305\pm0{,}032$\\
Ангармоническая поправка, $W<0{,}35$ & $+0{,}291\pm0{,}069$ & $+0{,}307\pm0{,}028$\\
\hline
\end{tabular}
\end{table}

В исходной гауссовской серии при $N=16$ на каждый класс
приходилось четыре блока по 1024 пробы. После проверки весов две выборки
при $k=4$ заменили четырьмя блоками по 8192 пробы на класс. Первые оценки ангармонической поправки к
$\delta f_{16}$ в двух сериях при $W<0{,}40$ составили $+0{,}640\pm0{,}090$
и $-0{,}233\pm0{,}263$~мэВ/атом. Разложение расхождения по вкладам заполнений
указало на четыре класса. В принятой нумерации узлов наборы позиций Al для
их представителей равны $\{0,13\}$, $\{0,2,13\}$, $\{0,1,14\}$ и
$\{0,2,13,15\}$. Для каждого из четырёх заново вычислили по 32\,768 проб
(четыре независимых блока по 8192) в каждой из двух серий. Ими заменили
прежние выборки этих классов. Для остальных 39 сохранили данные
соответствующей серии. Обе границы смещений оценивали на одних и тех же
пробах внутри каждой серии.

Стандартную ошибку итоговой величины $\theta$ оценивали блочным
джекнайфом. Поочерёдно исключали один блок $b$ у каждого представителя
и повторяли нормировку весов, логарифмирование, сумму по химическим
заполнениям и взятие разностей. Если $\theta_{(-b)}$~--- полученное
значение, $\overline{\theta}_{(-)}$~--- среднее по исключениям, то при
$B_{\rm bl}$ блоках
\[
 {\rm SE}^2=\frac{B_{\rm bl}-1}{B_{\rm bl}}
 \sum_{b=1}^{B_{\rm bl}}(\theta_{(-b)}-\overline{\theta}_{(-)})^2 .
\]
Так сохраняется корреляция между гармоникой и полным MTP, оценёнными на
одних смещениях. Даже после уточнения серия со Стьюдентом при $k=2$ и
$W<0{,}40$ дала эффективное число проб для общей химической и координатной
суммы около 565, с весом одной пробы до $3{,}96\%$. Для $k=3$ и 4
в этой серии оно превысило 7300. Обе независимые оценки дают около
$+0{,}30$~мэВ/атом, и этот сдвиг сохраняется при сужении допустимой
области (табл.~\ref{tab:small16}). Он характеризует
форму $F(x)$ для трёх заданных составов в ячейке из 16 атомов. Непрерывное
положение двух контактов при $p=0$ определяет отдельный расчёт с флуктуациями
состава, тепловых смещений и объёма.
